% !TeX root = ../../Book4.tex
% 纲要标题：### 12.1 矩映射的局部结构与 Morse–Bott 工具
% 纲要位置：几何学计划纲要.md，第 3919 行。
% 正文状态：未撰写；本文件目前是框架。

% 纲要内容开始
% 
% 就地证明凸性论证需要的局部标准形和次水平集变化，引入第二卷普通 Morse 工具时明确新增的临界子流形与负法丛。
% 
% * 固定点附近的权模型与相容复法丛
% * 正则轨道与局部辛标准形
% * 临界子流形、法向 Hessian 与 Morse–Bott 引理
% * 负法丛盘丛与跨临界值的次水平集变化
% * 一般圆方向矩映射分量的 Morse–Bott 性
% * 指标和余指标的偶性及其拓扑后果
% * 无指标一和无余指标一情形的连通水平集论证
% 
%   主范围为紧连通 Hamilton 环面作用的矩映射分量；核验极值分量、连通性与偶指标条件，不把连通水平集结论宣称为任意 Morse–Bott 函数都成立。
% 
% 纲要内容结束

\section{矩映射的局部结构与 Morse–Bott 工具}
\label{g4:sec:1201}

待撰写

\subsection{固定点附近的权模型与相容复法丛}
\label{g4:subsec:1201:01}
待撰写

\subsection{正则轨道与局部辛标准形}
\label{g4:subsec:1201:02}
待撰写

\subsection{临界子流形、法向 Hessian 与 Morse–Bott 引理}
\label{g4:subsec:1201:03}
待撰写

\subsection{负法丛盘丛与跨临界值的次水平集变化}
\label{g4:subsec:1201:04}
待撰写

\subsection{一般圆方向矩映射分量的 Morse–Bott 性}
\label{g4:subsec:1201:05}
待撰写

\subsection{指标和余指标的偶性及其拓扑后果}
\label{g4:subsec:1201:06}
待撰写

\subsection{无指标一和无余指标一情形的连通水平集论证}
\label{g4:subsec:1201:07}
待撰写
% 纲要条目排为编号小节标题，并进入目录与书签；说明段仍保留于纲要注释中。
